\subchapter{Board setup}{Objective: setup your board to be used during all
following labs}

During this lab, you will:
\begin{itemize}
  \item Build a Yocto image and the associated SDK
	\item Deploy it on the board
	\item Setup SSH connection with the board
\end{itemize}

\section{Build Yocto}

Go to the \code{$HOME/__SESSION_NAME__-labs/board_setup} directory.

Install the kas command, either from you distribution repositories or from
\href{pypi.org}{https://pypi.org/project/kas/}

\begin{bashinput}
$ python3 -m venv .venv
$ . ./.venv/bin/activate
$ pip install kas
\end{bashinput}

Then run kas to build the Yocto image. This might take a few hours.

\begin{bashinput}
$ kas build security-training-imx93-frdm.yaml
\end{bashinput}

\section{Set up the SD card}

We will now use an SD card to store the bootloader, kernel and root filesystem
files.
The SD card image has been generated in the build folder and is named
\code{security-training-image-freiheit93.rootfs.wic}.

Now flash the image with the following command:
\begin{bashinput}
sudo dd if=build/tmp-glibc/deploy/images/freiheit93/security-training-image-freiheit93.rootfs.wic of=/dev/mmcblkX conv=fdatasync bs=1M status=progress
\end{bashinput}

\section{Setting up serial communication with the board}
To set up serial communication with the board, simply connect the USB-C cable to
the port labeled DEBUG.

Once the cable is plugged in, a new serial port should appear:
\code{/dev/ttyACM0}.
You can also see this device appear by looking at the output of \code{dmesg}.

To communicate with the board through the serial port, install a
serial communication program, such as \code{picocom}:

\begin{bashinput}
sudo apt install picocom
\end{bashinput}

If you run \code{ls -l /dev/ttyACM0}, you can also see that only
\code{root} and users belonging to the \code{dialout} group have
read and write access to this file.
Therefore, you need to add your user to the \code{dialout} group:

\begin{bashinput}
sudo adduser $USER dialout
\end{bashinput}

{\bf Important}: for the group change to be effective, in Ubuntu 18.04, you have to
{\em completely reboot} the system \footnote{As explained on
\url{https://askubuntu.com/questions/1045993/after-adding-a-group-logoutlogin-is-not-enough-in-18-04/}.}.
A workaround is to run \code{newgrp dialout}, but it is not global.
You have to run it in each terminal.

Now, you can run \code{picocom -b 115200 /dev/ttyACM0}, to start serial
communication on \code{/dev/ttyACM0}, with a baudrate of \code{115200}.
If you wish to exit \code{picocom}, press \code{[Ctrl][a]} followed by
\code{[Ctrl][x]}.

There should be nothing on the serial line so far, as the board is not
powered up yet.

\section{Boot}
Insert the SD card in the dedicated slot on the imx93-frdm.
Make sure that switches 1 and 2 in the center of the board are set to ON, and
switches 3 and 4 are set to OFF.

Plug in the power USB-C.
You should see boot messages on the console.

During the very first boot, a full re-labelling of SELinux file context will
occur, followed by a reboot.
The whole process should take about 2 minutes.

Wait until the login prompt, then enter \code{root} as user.
Congratulations! The board has booted and you now have a shell.

\section{SSH access}

Plug a network cable on any of the two Ethernet ports of the board.
If the network is providing DHCP, the board should automatically get an IP
address; otherwise you will need to set it up by manually.

You can run the following command on the board to obtain the board IP address

\begin{bashinput}
# ip addr show eth0 | grep "inet "
\end{bashinput}

On your computer, you should be able to connect to the board using SSH:

\begin{bashinput}
$ BOARD=root@192.168.10.89
$ ssh root@$BOARD
root@freiheit93:~# 
\end{bashinput}

\section{Installing SDK}

Yocto has also been generating an SDK, providing a toolchain that can be used to
cross-compile software for your target.
As we will need to compile code during various labs, you should install this SDK
now:

\begin{bashinput}
$ /build/tmp-glibc/deploy/sdk/oecore-security-training-image-x86_64-cortexa55-freiheit93-toolchain-1.0.sh -d $HOME/__SESSION_NAME__-labs/sdk/
\end{bashinput}

You can then use the SDK by sourcing the environment file:

\begin{bashinput}
$ . $HOME/__SESSION_NAME__-labs/sdk/environment-setup-cortexa55-oe-linux
\end{bashinput}
